\documentclass[../main.tex]{subfiles}

\begin{document}
Prices in an open market reflects all of the available information regarding assets exchanged in an economy \cite{Malkiel2003}. When new information becomes available, all actors in the economy update their positions and prices adjust accordingly, which makes beating the markets consistently impossible. However, the definition of "new information" might change as new information retrieval technologies become available and early-adoption of such technologies might provide an advantage in the short-term. 

Analysis of financial texts, be it news, analyst reports or official company announcements is a possible source of new information. With unprecedented amount of such text being created every day, manually analyzing these and deriving actionable insights from them is too big of a task for any single entity. Hence, automated sentiment or polarity analysis of texts produced by financial actors using natural language processing (NLP) methods has gained popularity during the last decade \cite{Guo2016}.

The principal research interest for this thesis is the polarity analysis, which is classifying text as positive, negative or neutral, in a specific domain. It requires to address two challenges: 1) The most sophisticated classification methods that make use of neural nets require vast amounts of labeled data and labeling financial text snippets requires costly expertise. 2) The sentiment analysis models trained on general corpora are not suited to the task, because financial texts have a specialized language with unique vocabulary and have a tendency to use vague expressions instead of easily-identified negative/positive words. 

Using carefully crafted financial sentiment lexicons such as Loughran and McDonald (2011) \cite{Loughran2011} may seem a solution because they incorporate existing financial knowledge into textual analysis. However, they are based on "word counting" methods, which come short in analyzing deeper semantic meaning of a given text.

NLP transfer learning methods look like a promising solution to both of the challenges mentioned above, and are the focus of this thesis. The core idea behind these models is that by training language models on very large corpora and then initializing down-stream models with the weights learned from the language modeling task, a much better performance can be achieved. The initialized layers can range from the single word embedding layer \cite{Peters2018} to the whole model \cite{Howard2018}. This approach should, in theory, be an answer to the scarcity of labeled data problem. Language models don't require any labels, since the task is predicting the next word. They can learn how to represent the semantic information. That leaves the fine-tuning on labeled data only the task of learning how to use this semantic information to predict the labels.

One particular component of the transfer learning methods is the ability to further pre-train the language models on domain specific unlabeled corpus. Thus, the model can learn the semantic relations in the text of the target domain, which is likely to have a different distribution than a general corpus. This approach is especially promising for a niche domain like finance, since the language and vocabulary used is dramatically different than a general one.   

The goal of this thesis is to test these hypothesized advantages of using and fine-tuning pre-trained language models for financial domain. For that, sentiment of a sentence from a financial news article towards the financial actor depicted in the sentence will be tried to be predicted, using the Financial PhraseBank created by Malo et al. (2014) \cite{Malo2014} and FiQA Task 1 sentiment scoring dataset \cite{Maia2018}. 

The main contributions of this thesis are the following:
\begin{itemize}
    \item We introduce FinBERT, which is a language model based on BERT for financial NLP tasks. We evaluate FinBERT on two financial sentiment analysis datasets. 
    \item We achieve the state-of-the-art on FiQA sentiment scoring and Financial PhraseBank.
    \item We implement two other pre-trained language models, ULMFit and ELMo for financial sentiment analysis and compare these with FinBERT.  
    \item We conduct experiments to investigate several aspects of the model, including: effects of further pre-training on financial corpus, training strategies to prevent catastrophic forgetting and fine-tuning only a small subset of model layers for decreasing training time without a significant drop in performance. 
\end{itemize}

The rest of the thesis is structured as follows: First, relevant literature in both financial polarity analysis and pre-trained language models are discussed (Section \ref{sec:related_lit}). Then, the evaluated models are described (Section \ref{sec:method}). This is followed by the description of the experimental setup being used (Section \ref{exp_setup}). In Section \ref{exp_results}, we present the experimental results on the financial sentiment datasets. Then we further analyze FinBERT from different perspectives in Section \ref{exp_analysis}. Finally, we conclude with Section \ref{conclusion}.

\end{document}